% Generated draft: requires mathematical and source review.
\documentclass[11pt]{article}
\usepackage[margin=25mm]{geometry}
\usepackage{fontspec}
\setmainfont{Latin Modern Roman}
\usepackage{amsmath,amssymb,mathtools,unicode-math}
\setmathfont{Latin Modern Math}
\usepackage{enumitem,needspace,xurl,hyperref}
\setlength{\parindent}{0pt}
\setlength{\parskip}{5pt}
\setlength{\emergencystretch}{3em}
\newcommand{\N}{\mathbb{N}}
\newcommand{\Z}{\mathbb{Z}}
\newcommand{\Q}{\mathbb{Q}}
\newcommand{\R}{\mathbb{R}}
\newcommand{\C}{\mathbb{C}}
\newcommand{\F}{\mathbb{F}}
\newcommand{\A}{\mathbb{A}}
\newcommand{\PP}{\mathbb{P}}
\newcommand{\E}{\mathbb{E}}
\newcommand{\eps}{\varepsilon}
\newcommand{\dd}{\,\mathrm{d}}
\begin{document}
\section*{UnsolvedMath ID: 11000091}
\textbf{Source problem code:} catalogue entry
\subsection{Definitions and mathematical statement}
The target is to compute the graded invariant printed as
\[
 H^\bullet_\infty(T_3^c).
\]
Its coefficient system and the meaning of the infinity decoration are not defined in the extract. In particular, cohomology at infinity, compactly supported cohomology and locally finite homology are different theories and must not be silently interchanged. If the intended theory is cohomology at infinity, its usual meaning is \(H^k_\infty(X;A)=\varinjlim_{K\subset X\ \mathrm{compact}}H^k(X\setminus K;A)\), with the specified coefficient group \(A\); this is only a conditional interpretation.

In the standard geometric setting, \(T_g^c\) is the Torelli space of stable curves of compact type, with a symplectic integral homology marking, and \(T_g\) is its smooth-curve locus. Compact type means that the dual graph is a tree. The source relates their homology by a Gysin sequence and says its use requires understanding the topology of the reducible divisor \(T_g^{c,\mathrm{red}}\). This divisor is built from products of lower-genus compact-type Torelli spaces, with combinatorics described by
\[
 K_{\mathrm{sep}}(S)/I_g.
\]
Here \(K_{\mathrm{sep}}(S)\) is interpreted as the complex of separating curves, with disjointness simplices, and \(I_g\) as the Torelli group acting on it. The extracted final quotient uses an undecorated \(T_g\); the group symbol \(I_g\) merely distinguishes that role from the space, subject to checking the lost typography. The precise marking, strata/gluing, Gysin and homology-at-infinity conventions are absent. Computing the genus-three invariant is the question; no divisor decomposition or result is reconstructed beyond the supplied description.
\subsection{Short English statement}
Compute the source's infinity-decorated graded invariant of compact-type genus-three Torelli space, using the reducible divisor and separating-curve combinatorics once the missing homology conventions are fixed.
\subsection{Sources}
\begin{itemize}
\item \emph{Problems on Mapping Class Groups and Related Topics}, edited by Farb, Problem 4.6, supplied PDF page 77. \url{https://www.math.uchicago.edu/~farb/papers/mcgbook.pdf}. Recorded access 2026-07-29; extraction pdf-font-structure. This page identifies the occurrence despite repeated local numbering. No individual chapter author or separate reference is supplied in the record.
\item Local UnsolvedMath registry, \texttt{data/unsolvedmath/problems.json}, permanent ID \(11000091\), exact source code \texttt{catalogue entry}; supplied OPEN-TRIAGE summary reports unverified status.
\end{itemize}
\end{document}
